% econometrics.tex — Structural Econometric Specification
% Fee Concentration Risk and LP Exit: Quadratic Deviation Model
% Source: lp-insurance-demand.pdf (2026-03-05)

\section{Methodology and Data}\label{sec:econometric-full}

This section presents the full econometric identification of insurance demand
from fee concentration risk.

The primary finding is that, for the ETH/USDC 30bps pool we study, there
exists a concentration threshold ($\DeltaStar \approx 0.09$) beyond which
JIT arrival crowds out passive liquidity providers. In pools where trading
volume is inelastic to liquidity depth \cite{capponi2024jit}, JIT positions
dilute passive LPs' effective fee rate without generating proportional
volume---shifting concentration from a signal of pool health to a driver of
LP exit. This regime shift is the empirical foundation of insurance demand.

The goal is to identify precisely where this transition occurs and to
translate the excess exit risk into a premium for fee concentration
derivatives. This motivates two questions:

\begin{enumerate}
  \item Does fee concentration risk ($\AT$) drive passive LP exits in
    Uniswap V3?
  \item At what concentration threshold does the exit effect activate, and
    can it be priced as an insurance premium?
\end{enumerate}

\subsubsection{Unit of Observation}

\textbf{Position-day panel.} One observation per position per day over a
41-day window (2025-12-05 to 2026-01-14). Binary outcome:
$\text{exit}_{i,t} \in \{0,1\}$. Treatment: fee concentration deviation from
the competitive null, measured at the pool-day level with lag.

\subsubsection{Outcome Variable}

Binary exit indicator:
\begin{equation}\label{eq:exit-indicator}
  Y_{i,t} = \begin{cases}
    1 & \text{if position } i \text{ burns on day } t \\
    0 & \text{otherwise}
  \end{cases}
\end{equation}

Panel size: 3{,}365 position-day observations, 597 exit events, 40 cluster
days, from 600 positions in the ETH/USDC 30bps pool.

%% ─── 2. The Fee Concentration Index ─────────────────────────────────────────

\subsection{Fee Concentration Index (Econometric Definition)}\label{subsec:econ-fci}

\begin{definition}[Fee Concentration Index]\label{def:econ-fci}
For a pool with $\PosCount$ active positions on day $t$, the fee concentration
index is:
\begin{equation}\label{eq:econ-at}
  \AT = \sqrt{\sum_{k=1}^{\PosCount} \thetapos_k \, x_k^2}
\end{equation}
where $x_k = \text{fee}_k / \sum_j \text{fee}_j$ is position $k$'s fee share
and $\thetapos_k = 1/\blocklife_k$ is its turnover rate.
\end{definition}

$\AT$ aggregates two dimensions of competitive pressure: \emph{who earns the
fees} ($x_k$) and \emph{how fast positions turn over} ($\thetapos_k$). High
$\AT$ means a few short-lived positions capture most fee revenue.

\subsubsection{The Equal-Share Null}

\begin{definition}[Ma--Crapis Null]\label{def:econ-null}
Under the Ma \& Crapis (2024) symmetric equilibrium with $\PosCount$ identical
LPs, each earns equal fee share $x_k = 1/\PosCount$:
\begin{equation}\label{eq:econ-null}
  \ATnull = \sqrt{\frac{\sum_k \thetapos_k}{\PosCount^2}}
\end{equation}
This is the competitive baseline---the value $\AT$ takes when fee
concentration is zero.
\end{definition}

\subsubsection{Concentration Deviation}

\begin{definition}[Concentration Deviation]\label{def:econ-deviation}
The treatment variable is the excess concentration above the null:
\begin{equation}\label{eq:econ-deviation}
  \Delta_t \equiv A_{T,t}^{\text{real}} - A_{T,t}^{1/N}
\end{equation}
computed from the \textbf{same} set of positions on each day.
$\Delta_t > 0$ indicates concentration above the competitive equilibrium;
$\Delta_t < 0$ indicates more equal fee distribution than the symmetric
benchmark.
\end{definition}

\subsubsection{Empirical Properties of $\Delta_t$}

Data source: Dune query 6783604, joining \texttt{Collect}, \texttt{Burn}, and
\texttt{Mint} events for the ETH/USDC 30bps pool
(\texttt{0x8ad599c3\ldots eb6e6d8}). Fee amounts isolated by subtracting Burn
principal from Collect amounts. Token ordering: token0 = USDC (6 dec),
token1 = WETH (18 dec).

\begin{center}
\begin{tabular}{@{}ll@{}}
\toprule
Statistic & Value \\
\midrule
Days with $\Delta > 0$ & 26/41 (63\%) \\
Mean $A_T^{\text{real}} / A_T^{1/N}$ & $2.65\times$ \\
Median deviation & $+0.0004$ \\
Range & $[-0.032,\; +0.158]$ \\
Positions per day & 13--65 \\
\bottomrule
\end{tabular}
\end{center}

The real $\AT$ exceeds the null on most days, confirming empirically that fee
shares are asymmetric---consistent with Aquilina et al.\ finding that 7\% of
LP addresses capture 80\% of fees.

%% ─── 3. Theoretical Foundations ─────────────────────────────────────────────

\subsection{Theoretical Foundations}\label{subsec:econ-theory}

Five papers independently validate the components of $\AT$ and the exit
mechanism.

\subsubsection{Capponi \& Zhu (2024): Optimal Exit Threshold}

\emph{``Optimal Exiting for Liquidity Provision in Constant Function Market Makers.''}

The LP solves an optimal stopping problem with exit thresholds $p_U^*$ and
$p_L^*$. The key comparative static:
\begin{equation}\label{eq:capponi-exit}
  \frac{\partial p_U^*}{\partial \phi} > 0
\end{equation}
A higher effective fee rate $\hat{\phi}$ raises the exit threshold (LP stays
longer). When $\AT$ rises, passive LP's effective fee rate
$\hat{\phi}_i = x_i \cdot \phi_{\text{pool}}$ drops because $x_i < 1/\PosCount$.
Lower $\hat{\phi} \Rightarrow$ lower $p_U^* \Rightarrow$ earlier exit.

\textbf{Implication for $\AT$:} $\AT$ enters the exit decision through
$\hat{\phi}$. This predicts $\beta_1 > 0$ for sufficiently high concentration.

\subsubsection{Capponi, Jia \& Zhu (2024): JIT Fee Dilution}

\emph{``The Paradox of Just-in-Time Liquidity in Decentralized Exchanges.''}

JIT liquidity creates fee share asymmetry: the passive LP's share is
$d_P/(d_P + d_J)$ where $d_J$ is JIT depth. This is exactly the $x_k$
component of $\AT$. JIT entry mechanically raises $\AT$ by concentrating fees
in short-lived positions.

\textbf{Implication for identification:} JIT activity is a natural instrument
for $\AT$. Our JIT control analysis (32 JIT positions, 5.3\% share) shows JIT
count has an independent positive effect on exits
($\beta_{\text{JIT}} = 0.177$, $p = 0.031$), but does not explain the $\AT$
effect---confirming that $\AT$ captures concentration \emph{beyond} JIT.

\subsubsection{Ma \& Crapis (2024): The Competitive Null}

\emph{``Decentralized Exchange Design: Permissionless Liquidity Provision.''}

$\PosCount$ symmetric LPs in free-entry equilibrium: each earns
$\Pi_i = (1/\PosCount^2) f B_D$. The equal-share outcome $x_k = 1/\PosCount$
is the \textbf{null hypothesis} of $\AT$. Our deviation
$\Delta_t = A_T^{\text{real}} - A_T^{1/N}$ tests departure from this
equilibrium directly.

\textbf{Implication for the model:} $\Delta_t = 0$ means the pool operates at
the competitive equilibrium. $\Delta_t > 0$ means fee concentration exceeds
what symmetric competition predicts. The deviation is the correct treatment
variable---not $\AT$ in levels.

\subsubsection{Aquilina, Foley, Gambacorta \& Krekel (2024): Empirical Validation}

\emph{``Decentralised Dealers? Examining Liquidity Provision in Decentralised
Exchanges.''} BIS Working Paper No.~1227.

Sophisticated LPs (7\% of addresses) provide 65--85\% of liquidity, use tick
ranges $3.4\times$ tighter, hold positions $7.5\times$ shorter, and capture
75--90\% of fees. This validates that $x_k$ in $\AT$ is highly skewed, not
symmetric.

Their profitability decomposition
$R_{\text{net}} = \text{FeeYield} + \text{IL} - \text{Gas}$ shows that passive
LPs earn negative returns when sophisticated LPs dominate---precisely when
$\Delta_t$ is large.

\subsubsection{Bichuch \& Feinstein (2024): Fee-as-Derivative Pricing}

\emph{``The Price of Liquidity: Implied Volatility of AMM Fees.''}

The LP value function is an optimal stopping problem with implied fee rate
$F(p_x, p_y) = p_y \cdot \ell(p_x/p_y)$ where $\ell$ is instantaneous LVR.
Theorem~4.2: a risk-neutral LP is indifferent across all stopping times. This
means \emph{any} deterministic exit driver must come from fee rate
heterogeneity across LPs---exactly what $\AT$ measures.

Their fixed-for-floating fee swap (Corollary~5.5) is structurally identical to
ThetaSwap's insurance product: the LP exchanges volatile realized fees for a
fixed premium.

\textbf{Implication:} $\AT$ measures the \emph{distribution} of the pool fee
rate across LPs. Bichuch \& Feinstein's framework prices the pool-level fee;
$\AT$ determines who receives it.

%% ─── 4. Economic Model ─────────────────────────────────────────────────────

\subsection{Economic Model}\label{subsec:econ-model}

\subsubsection{Environment and Actors}

Single Uniswap V3 pool (ETH/USDC 30bps). Two agent types:

\begin{enumerate}
  \item \textbf{Passive LP (PLP):} Commits liquidity for extended periods.
    Observes pool state (reserves, price, accumulated fees). Does \emph{not}
    observe $\AT$ directly but reacts to its symptoms (lower fee revenue).
    Decision: when to exit.

  \item \textbf{JIT/Sophisticated LP:} Not explicitly modeled as strategic.
    Enters as a latent force that raises $\AT$ by capturing disproportionate
    fee shares. Characterized by short blocklife and tight tick ranges
    (Aquilina et al.).
\end{enumerate}

\subsubsection{Equilibrium Concept}

Competitive (price-taking) per Capponi \& Zhu (2024). The PLP takes fee rate,
price process, and concentration risk as given and solves an optimal stopping
problem.

\textbf{Assumption} (Price-Taking LP). \emph{The PLP solves:}
\begin{equation}\label{eq:plp-stopping}
  \max_{\tau}\; \mathbb{E}\!\left[
    \int_0^{\tau} e^{-rt} \hat{\phi}_i \, dF_t - dL_t
    \;\middle|\; \mathcal{F}_0
  \right]
\end{equation}
where $\hat{\phi}_i = x_i \cdot \phi_{\text{pool}}$ is the LP's effective fee
rate given fee share $x_i$, $F_t$ is cumulative fee revenue, and $L_t$ is
impermanent loss.

%% ─── 5. Stochastic Model ───────────────────────────────────────────────────

\subsection{Stochastic Model}\label{subsec:econ-stochastic}

\subsubsection{Unobserved Heterogeneity ($\eta$)}

\begin{enumerate}
  \item \textbf{LP risk aversion} ($\eta_i^{\text{risk}}$): heterogeneous
    preferences. Some PLPs tolerate more concentration. NFT Position Manager
    masks identity.

  \item \textbf{Outside options} ($\eta_i^{\text{outside}}$): competing pool
    returns, staking yields, CEX opportunities.
\end{enumerate}

Bias direction: \textbf{toward zero} (conservative). Risk-averse LPs avoid
entering during high-$\AT$ periods \emph{and} exit faster, attenuating
$\beta$.

\subsubsection{Agent Uncertainty}

Future price process only ($u_t \sim$ GBM, per Capponi \& Zhu). Drives IL
realization.

\subsubsection{Implied Error Structure}

\begin{equation}\label{eq:error-structure}
  \varepsilon_i =
    \underbrace{\eta_i^{\text{risk}} + \eta_i^{\text{outside}}}_{\text{unobserved heterogeneity}}
    + \underbrace{\nu_i}_{\text{measurement error}}
\end{equation}

%% ─── 6. Estimation Strategy ────────────────────────────────────────────────

\subsection{Estimation Strategy}\label{subsec:econ-estimation}

\subsubsection{Model Progression}

We estimate three nested specifications to motivate the quadratic model.

\paragraph{Specification 1: Linear (Point Lag).}
\begin{equation}\label{eq:spec-linear-point}
  P(\text{exit}_{i,t} = 1) = \sigma\bigl(
    \beta_0 + \beta_1 A_{T,t-\ell} + \beta_2 \,\text{IL}_t
    + \beta_3 \log(\text{age}_{i,t})
  \bigr)
\end{equation}

\textbf{Result 1} (Linear A\_T). $\beta_1 = -3.93$, cluster $p = 0.072$.
Marginally significant but \textbf{wrong sign}---higher concentration
associated with fewer exits. Robust across lags 1--7 (all negative, 4/5
significant).

\paragraph{Specification 2: Linear Deviation from $1/\PosCount$.}
\begin{equation}\label{eq:spec-linear-dev}
  P(\text{exit}_{i,t} = 1) = \sigma\bigl(
    \beta_0 + \beta_1 \Delta_{t-\ell} + \beta_2 \,\text{IL}_t
    + \beta_3 \log(\text{age}_{i,t})
  \bigr)
\end{equation}
where $\Delta_t = A_T^{\text{real}} - A_T^{1/N}$ from the same position set.

\textbf{Result 2} (Linear Deviation). $\beta_1 = -5.44$, cluster $p = 0.109$.
Still negative. But \textbf{dose-response analysis reveals non-monotonicity}:

\begin{center}
\begin{tabular}{@{}llll@{}}
\toprule
Quartile & Mean $\Delta$ & Exit Rate & Pattern \\
\midrule
Q1 (below null) & $-0.009$ & 17.75\% & baseline \\
Q2 (at null) & $-0.000$ & 17.87\% & baseline \\
\textbf{Q3 (mild conc.)} & $\mathbf{+0.002}$ & \textbf{20.56\%} & \textbf{exits increase} \\
Q4 (extreme conc.) & $+0.025$ & 14.37\% & exits decrease \\
\bottomrule
\end{tabular}
\end{center}

Q3 shows the predicted Capponi \& Zhu pattern. Q4 reverses---motivating the
quadratic specification.

\paragraph{Specification 3: Quadratic Deviation (Primary Model).}
\begin{equation}\label{eq:quadratic-hazard-full}
  \boxed{P(\text{exit}_{i,t} = 1) = \sigma\bigl(
    \beta_0 + \beta_1 \Delta_{t-\ell} + \beta_2 \Delta_{t-\ell}^2
    + \beta_3 \,\text{IL}_t + \beta_4 \log(\text{age}_{i,t})
  \bigr)}
\end{equation}

This is the \textbf{primary specification}. The quadratic term captures the
inverted-U: shelter at low concentration, Capponi exit at high concentration.

\subsubsection{Estimation Method}

Logit MLE via Newton--Raphson (IRLS) with \textbf{day-clustered sandwich
standard errors} (41 clusters):
\begin{equation}\label{eq:clustered-var}
  \widehat{\text{Var}}(\hat{\beta}) = H^{-1}
    \left(\sum_{g=1}^{G} S_g S_g'\right)
    H^{-1} \cdot \frac{G}{G-1} \cdot \frac{n}{n-k}
\end{equation}
where $S_g = \sum_{i \in g} (y_i - \hat{p}_i) x_i$ is the cluster-level score
and $H$ is the Hessian.

Implemented in JAX 0.9.1 CPU. Code: \texttt{econometrics/hazard.py}.

%% ─── 7. Results ─────────────────────────────────────────────────────────────

\subsection{Results}\label{subsec:econ-results}

\subsubsection{Main Result: Quadratic Deviation Model}

\begin{center}
\begin{tabular}{@{}cccccl@{}}
\toprule
Lag & $\hat{\beta}_1$ (linear) & $p$ & $\hat{\beta}_2$ (quadratic) & $p$ & Turning point \\
\midrule
\textbf{1} & $\mathbf{-23.18}$ & $\mathbf{0.012}$ & $\mathbf{+129.20}$ & $\mathbf{0.030}$ & $\mathbf{0.090}$ \\
\textit{2} & \textit{$-43.42$} & \textit{0.016} & \textit{$+226.92$} & \textit{0.065} & \textit{0.096} \\
\textbf{3} & $\mathbf{-32.44}$ & $\mathbf{0.001}$ & $\mathbf{+205.34}$ & $\mathbf{0.004}$ & $\mathbf{0.079}$ \\
5 & $+3.81$ & 0.830 & $-76.31$ & 0.470 & --- \\
7 & $-5.55$ & 0.730 & $-16.46$ & 0.900 & --- \\
\bottomrule
\end{tabular}
\end{center}

Both terms of the quadratic deviation model are statistically significant at
the economically relevant lag horizon (1--3 days). The turning point clusters
at $\DeltaStar \approx 0.08$--$0.10$. The signal dissipates at lags 5--7,
consistent with PLPs responding to recent concentration shocks, not stale ones.

\subsubsection{Full Lag-1 Coefficients}

\begin{center}
\begin{tabular}{@{}lcrl@{}}
\toprule
Parameter & Estimate & Cluster SE & Interpretation \\
\midrule
$\beta_0$ (intercept) & varies & --- & baseline exit probability \\
$\beta_1$ ($\Delta$) & $-23.18$ & 9.21 & shelter below turning point \\
$\beta_2$ ($\Delta^2$) & $+129.20$ & 59.69 & exit above turning point \\
$\beta_3$ (IL) & $+13.52$ & --- & higher IL $\to$ more exits \\
$\beta_4$ (log age) & $-0.42$ & --- & older positions less likely to exit \\
\bottomrule
\end{tabular}
\end{center}

$n = 3{,}365$; exits $= 597$; clusters $= 40$; pseudo-$R^2 = 0.350$;
mean $P(\text{exit}) = 0.177$.

%% ─── 8. Economic Interpretation ─────────────────────────────────────────────

\subsection{Economic Interpretation}\label{subsec:econ-interpretation}

\subsubsection{The Inverted-U: Two Regimes}

The quadratic model identifies two regimes separated by the turning point
$\DeltaStar \approx 0.09$:

\begin{enumerate}
  \item \textbf{Shelter regime} ($\Delta < \DeltaStar$): Mild fee concentration
    correlates with high trading volume. All LPs benefit from elevated fee
    revenue, even if shares are unequal. The marginal effect
    $\partial P / \partial \Delta < 0$: more concentration, fewer exits.

  \item \textbf{Capponi regime} ($\Delta > \DeltaStar$): Extreme concentration
    means passive LPs' effective fee rate $\hat{\phi}_i$ drops below their exit
    threshold. Per Capponi \& Zhu:
    $\partial p_U^* / \partial \phi > 0$, so lower $\phi \Rightarrow$ lower
    $p_U^* \Rightarrow$ earlier exit. The marginal effect flips:
    $\partial P / \partial \Delta > 0$.
\end{enumerate}

\subsubsection{Marginal Effect at the Turning Point}

The marginal effect of the quadratic model is:
\begin{equation}\label{eq:marginal-effect}
  \frac{\partial P(\text{exit})}{\partial \Delta}
    = (\beta_1 + 2\beta_2 \Delta) \cdot \bar{P}(1 - \bar{P})
\end{equation}

At the turning point $\DeltaStar = -\beta_1/(2\beta_2) = 0.090$:
\[
  \left.\frac{\partial P}{\partial \Delta}\right|_{\DeltaStar} = 0
\]

For $\Delta = 0.15$ (well into the Capponi regime):
\begin{align}
  \beta_1 + 2\beta_2(0.15) &= -23.18 + 2(129.20)(0.15) = 15.58
    \label{eq:marginal-015} \\
  \frac{\partial P}{\partial \Delta} &= 15.58 \times 0.177 \times 0.823
    = 2.27 \label{eq:marginal-value}
\end{align}

A 0.01 increase in $\Delta$ at this level raises exit probability by
${\sim}2.3$ percentage points.

\subsubsection{Insurance Pricing Translation}

For a passive LP experiencing $\Delta = 0.15$ (concentration 15\% above the
competitive null):

\begin{center}
\begin{tabular}{@{}ll@{}}
\toprule
Quantity & Value \\
\midrule
$\Delta P(\text{exit})$ per 0.01 $\Delta$ & $+2.3$ pp \\
Average remaining position hours & 48 \\
Hours of fee revenue at risk & $0.023 \times 48 = 1.1$ hrs \\
Fee revenue per hour (pool avg) & \$100 \\
Implied premium per position & \$110 \\
\bottomrule
\end{tabular}
\end{center}

This is the maximum willingness-to-pay for ThetaSwap's insurance: a PLP would
pay up to \$110 to eliminate the excess exit risk from concentration above the
turning point.

%% ─── 9. Specification Tests ─────────────────────────────────────────────────

\subsection{Specification Tests}\label{subsec:econ-tests}

\subsubsection{Test 1: Sign Restriction (Capponi \& Zhu Prediction)}

\textbf{Hypothesis 1.} $\beta_2 > 0$: the quadratic term is positive,
indicating that extreme concentration drives exits.

\textbf{Result: PASS.} $\hat{\beta}_2 = +129.20$, $p = 0.030$ at lag~1;
$+205.34$, $p = 0.004$ at lag~3. The Capponi exit mechanism is confirmed in
the high-concentration regime.

\subsubsection{Test 2: Dose-Response Monotonicity}

\textbf{Hypothesis 2.} Exit rates increase with $\Delta$ in the upper
quartiles.

\textbf{Result: PASS (inverted-U).} Q1--Q2 stable (17.8\%), Q3 peak (20.6\%),
Q4 drops (14.4\%). The non-monotonicity motivates the quadratic and is
captured by the model.

\subsubsection{Test 3: Lag Robustness}

\textbf{Hypothesis 3.} Both $\beta_1$ and $\beta_2$ are significant for at
least 3 of 5 lag windows.

\textbf{Result: PASS.} Lags 1, 2, 3 all have significant quadratic terms
($p < 0.07$). The turning point is stable at 0.08--0.10. Signal dissipates at
lags 5--7, consistent with short-memory PLP reaction.

\subsubsection{Test 4: IL Orthogonality}

\textbf{Hypothesis 4.} $\beta_1, \beta_2$ persist with and without the IL
control.

\textbf{Result: PASS.} IL coefficient ($\beta_3 = 13.52$) is positive
(higher IL $\to$ more exits) but does not absorb the concentration effect.
The quadratic structure is present in the $\AT$-only model as well.

\subsubsection{Test 5: JIT Confounding Control}

\textbf{Hypothesis 5.} The concentration effect is not an artifact of JIT
activity.

\textbf{Result: PASS.} Adding a JIT proxy (positions with blocklife $< 100$
blocks) as a control:
\begin{itemize}
  \item JIT count has independent positive effect on exits
    ($\beta_{\text{JIT}} = 0.177$, $p = 0.031$)
  \item The $\AT$ effect becomes \emph{more} significant
    ($p: 0.072 \to 0.022$) with JIT controlled
  \item Confirms: $\AT$ captures concentration \emph{beyond} JIT
\end{itemize}

\subsubsection{Test 6: Alternative Treatment Timing}

Three treatment constructions tested:

\begin{center}
\begin{tabular}{@{}llll@{}}
\toprule
Treatment & $\beta_1$ (linear) & $p$ & Finding \\
\midrule
Point A\_T (real, lag 1) & $-3.93$ & 0.072 & Negative (shelter) \\
Lifetime mean A\_T & $-7.60$ & 0.444 & Negative, insignificant \\
Deviation from $1/\PosCount$ & $-5.44$ & 0.109 & Negative (Q3 reveals non-linearity) \\
\textbf{Quadratic deviation} & \multicolumn{2}{l}{\textbf{See Result 3}} & \textbf{Inverted-U confirmed} \\
\bottomrule
\end{tabular}
\end{center}

%% ─── 10. Robustness: Five Linear Specifications ─────────────────────────────

\subsection{Robustness: Five Linear Specifications}\label{subsec:econ-robustness}

Before arriving at the quadratic model, we tested five linear specifications
that progressively refined the treatment variable. All failed to produce a
positive $\beta_1$, but each contributed to understanding the data-generating
process.

\subsubsection{Specification A: Point Real A\_T}

Treatment $= A_{T,t-\ell}^{\text{real}}$ (single day's concentration).

\begin{center}
\begin{tabular}{@{}ccccc@{}}
\toprule
Lag & $\hat{\beta}_1$ & Cluster SE & $p$ & Sig \\
\midrule
1 & $-3.928$ & 2.186 & 0.072 & * \\
2 & $-8.783$ & 2.227 & 0.0001 & *** \\
3 & $-0.414$ & 1.131 & 0.714 & \\
5 & $-6.523$ & 1.650 & 0.0001 & *** \\
7 & $-5.625$ & 2.261 & 0.013 & ** \\
\bottomrule
\end{tabular}
\end{center}

\textbf{Finding:} Consistently negative. A naive interpretation is that
concentration \emph{protects} LPs. But the proxy A\_T ($x_k = 1/\PosCount$)
showed the same pattern ($\beta_1 = -4.90$, $p = 0.15$), suggesting the level
of $\AT$ confounds pool health (volume) with concentration pressure.

\subsubsection{Specification B: Lifetime Mean A\_T}

Treatment $=$ mean $A_T^{\text{real}}$ over $[\text{mint}_i, t - \ell]$.
Motivated by the hypothesis that passive LPs react to accumulated exposure,
not point shocks.

\begin{center}
\begin{tabular}{@{}ccccc@{}}
\toprule
Lag & $\hat{\beta}_1$ & Cluster SE & $p$ & Sig \\
\midrule
1 & $-7.595$ & 9.913 & 0.444 & \\
3 & $-11.603$ & 7.959 & 0.145 & \\
5 & $-14.977$ & 8.021 & 0.062 & * \\
7 & $-3.042$ & 6.829 & 0.656 & \\
\bottomrule
\end{tabular}
\end{center}

\textbf{Finding:} Amplifies the negative effect but with much larger SEs
(reduced variation in running means). Marginal significance only at lag~5.

\subsubsection{Specification C: Deviation from $1/\PosCount$ (Wrong Null)}

First attempt at the deviation treatment, using the old proxy A\_T from a
\emph{different} position set (Q4 burns) as the null. The real and null A\_T
were computed from different populations---an apples-to-oranges comparison.

\textbf{Finding:} Real A\_T was $10$--$100\times$ smaller than the proxy null
(different position sets, different blocklives). Deviation was negative for
40/41 days. Only lag-1 showed positive $\beta_1$ ($+1.92$, $p = 0.43$).
\textbf{Discarded} after identifying the data mismatch.

\subsubsection{Specification D: Corrected Deviation (Same Position Set)}

Null A\_T recomputed from the \emph{same} Dune query (same positions, same
blocklives). Now 26/41 days have $\Delta > 0$, mean ratio $= 2.65\times$.

\textbf{Finding:} Still negative $\beta_1$ in linear model, but dose-response
revealed the Q3 peak (20.56\%) and Q4 dip (14.37\%) that motivated the
quadratic specification.

\subsubsection{Specification E: JIT-Controlled Linear}

Added daily JIT count (positions with blocklife $< 100$ blocks) as control.

\textbf{Finding:} $\beta_1$ became \emph{more} negative and more significant
($p: 0.072 \to 0.022$). JIT itself predicts exits but does not confound the
$\AT$ effect.

%% ─── 11. Literature Alignment ───────────────────────────────────────────────

\subsection{Literature Alignment}\label{subsec:econ-literature}

Each component of the quadratic deviation model is validated by at least one
paper:

\begin{center}
\begin{tabular}{@{}p{3.5cm}p{4cm}p{5.5cm}@{}}
\toprule
Paper & Component Validated & Connection to Our Result \\
\midrule
Capponi \& Zhu (2024) &
  Exit threshold: $\partial p_U^*/\partial\phi > 0$ &
  Confirmed in the Capponi regime ($\Delta > 0.09$):
  $\beta_2 > 0$ means extreme concentration drives exits \\[6pt]
Capponi, Jia \& Zhu (2024) &
  $x_k$ captures JIT fee dilution &
  JIT control shows $\AT$ effect is independent of JIT confounding;
  $x_k$ reflects broader concentration than just JIT \\[6pt]
Ma \& Crapis (2024) &
  $x_k = 1/\PosCount$ is the competitive null &
  Deviation $\Delta = A_T^{\text{real}} - A_T^{1/N}$ is the correct treatment;
  levels confound pool health \\[6pt]
Aquilina et al.\ (2024) &
  7\% of LPs capture 80\% of fees &
  Real $\AT$ is $2.65\times$ the null on average, confirming massive fee share
  asymmetry \\[6pt]
Bichuch \& Feinstein (2024) &
  Risk-neutral LP indifferent on stopping times &
  Any exit driver must come from fee \emph{heterogeneity} across LPs---exactly
  what $\Delta$ measures \\
\bottomrule
\end{tabular}
\end{center}

%% ─── 12. Connection to ThetaSwap ────────────────────────────────────────────

\subsection{Connection to ThetaSwap}\label{subsec:econ-thetaswap}

\begin{enumerate}
  \item \textbf{Insurance trigger:} The turning point $\DeltaStar \approx 0.09$
    defines when insurance becomes valuable. When concentration exceeds 9\%
    above the Ma--Crapis equilibrium, passive LP exit risk increases.
    ThetaSwap monitors $\Delta_t$ in real time and activates payouts when
    $\Delta > \DeltaStar$.

  \item \textbf{Premium calibration:} The quadratic marginal effect at any
    $\Delta$ translates directly to an insurance premium via
    \cref{eq:marginal-effect}. At $\Delta = 0.15$: \$110 per position per
    event.

  \item \textbf{Distinct from IL hedging:} The IL coefficient ($\beta_3$) is
    positive but does not absorb the concentration effect. $\AT$ captures a
    risk orthogonal to price-movement risk, confirming that ThetaSwap
    addresses a gap unmet by options or IL protection products.

  \item \textbf{Product structure:} Per Bichuch \& Feinstein's
    fixed-for-floating fee swap framework, ThetaSwap is a derivative that pays
    the PLP when realized fee share $x_i$ falls below $1/\PosCount$ by more
    than $\DeltaStar$. The premium is the implied volatility of fee shares.

  \item \textbf{Market sizing:} 63\% of days in our sample have
    $\Delta > 0$. Of these, days exceeding $\DeltaStar$ represent the
    insurance-triggering events. Aggregating the implied premium across 600
    positions gives the minimum viable insurance pool size.
\end{enumerate}

%% ─── 13. Data and Implementation ────────────────────────────────────────────

\subsection{Data and Implementation}\label{subsec:econ-data}

\begin{center}
\begin{tabular}{@{}llll@{}}
\toprule
Query & Description & Rows & Cost \\
\midrule
Q4v2 & Position-level burns, blocklife & 600 & 0.18 cr \\
Q5 & Daily IL proxy & 41 & 0.14 cr \\
Q6 & Per-position Collect fees + blocklife (real $\AT$ and null $\AT$ per day)
  & 41 (aggregated) & 3.3 cr \\
\bottomrule
\end{tabular}
\end{center}

Total cost: ${\sim}3.6$ credits.

\textbf{Pool:} ETH/USDC 30bps (\texttt{0x8ad599c3a0ff1de082011efddc58f1908eb6e6d8}).\\
\textbf{Window:} 2025-12-05 to 2026-01-14 (41 days).\\
\textbf{Implementation:} Python 3.11, JAX 0.9.1 CPU, \texttt{@functional-python}
patterns (frozen dataclasses, pure functions).
Code in \texttt{econometrics/\{types,data,ingest,hazard\}.py}. 52 unit tests
passing.

Key functions:
\begin{itemize}
  \item \texttt{build\_exit\_panel\_deviation()}: constructs position-day panel
    with $\Delta$ as treatment
  \item \texttt{logit\_mle\_quadratic()}: Newton--Raphson logit with
    $\Delta + \Delta^2$ and clustered SEs
  \item \texttt{marginal\_effect\_at\_means()}: translates $\beta$ to insurance
    premium
\end{itemize}

%% ─── 14. References ─────────────────────────────────────────────────────────

\subsection{Econometric References}

\begin{itemize}
  \item Aquilina, M., Foley, S., Gambacorta, L., \& Krekel, M.\ (2024).
    ``Decentralised Dealers? Examining Liquidity Provision in Decentralised
    Exchanges.'' \emph{BIS Working Paper} No.~1227.
  \item Bichuch, M.\ \& Feinstein, Z.\ (2024).
    ``The Price of Liquidity: Implied Volatility of AMM Fees.''
  \item Capponi, A.\ \& Zhu, B.\ (2024).
    ``Optimal Exiting for Liquidity Provision in Constant Function Market Makers.''
  \item Capponi, A., Jia, R., \& Zhu, B.\ (2024).
    ``The Paradox of Just-in-Time Liquidity in Decentralized Exchanges:
    More Providers Can Lead to Less Liquidity.''
  \item Ma, J.\ \& Crapis, H.\ (2024).
    ``Decentralized Exchange Design: Permissionless Liquidity Provision.''
  \item Reiss, P.C.\ \& Wolak, F.A.\ (2007).
    ``Structural Econometric Modeling: Rationales and Examples from
    Industrial Organization.'' \emph{Handbook of Econometrics}, Vol.~6A, Ch.~64.
\end{itemize}
